% Дробление шага: пробный шаг принимается, когда график phi(alpha) опускается
% ниже прямой условия дробления.
\begin{figure}[htbp]
\centering
\begin{tikzpicture}
\begin{axis}[
  width=10.2cm, height=5.8cm,
  axis lines=middle,
  xmin=-0.12, xmax=2.45, ymin=-0.45, ymax=2.45,
  xlabel={$\alpha$},
  xlabel style={anchor=north east},
  xtick={1.0, 2.0}, xticklabels={$\alpha_k = \alpha^{(0)}\lambda$, $\alpha^{(0)}$},
  ytick=\empty, clip=false, samples=80,
]
  % phi(alpha) = f(x_k - alpha grad f)
  \addplot[very thick, blue!60!black, domain=0:2.1] {1 - 2*x + 1.2*x^2};
  \node[font=\footnotesize, blue!60!black, anchor=west] at (axis cs:1.75,1.55)
    {$\varphi(\alpha) = f\big(x_k - \alpha \nabla f(x_k)\big)$};

  % прямая условия дробления
  \addplot[very thick, red!65!black, dashed, domain=0:2.3] {1 - 0.4*x};
  \node[font=\footnotesize, red!65!black, anchor=north west]
    at (axis cs:2.05,0.12)
    {$\varphi(0) - \eps\alpha\norm{\nabla f(x_k)}^2$};

  % касательная в нуле
  \addplot[gray, dotted, thick, domain=0:0.9] {1 - 2*x};
  \node[font=\scriptsize, gray, anchor=east] at (axis cs:0.88,-0.75) {};

  % отвергнутый шаг
  \draw[dashed, gray] (axis cs:2.0,0.2) -- (axis cs:2.0,1.8);
  \fill[red!70!black] (axis cs:2.0,1.8) circle (2.0pt);
  \node[font=\scriptsize, red!70!black, anchor=south] at (axis cs:2.0,1.88)
    {отвергнут};

  % принятый шаг
  \draw[dashed, gray] (axis cs:1.0,0.2) -- (axis cs:1.0,0.6);
  \fill[teal!70!black] (axis cs:1.0,0.2) circle (2.0pt);
  \node[font=\scriptsize, teal!70!black, anchor=south east]
    at (axis cs:0.97,0.28) {принят};
\end{axis}
\end{tikzpicture}
\caption{Дробление шага. Пробное значение $\alpha$ принимается, если график
$\varphi$ в этой точке лежит \emph{не выше} штриховой прямой. Прямая проходит
через $\varphi(0)$ с наклоном $-\eps\norm{\nabla f(x_k)}^2$, то есть она
положе касательной (наклон $-\norm{\nabla f(x_k)}^2$) ровно в $1/\eps$ раз —
поэтому вблизи нуля условие заведомо выполняется, и дробление конечно.}
\label{fig:step-halving}
\end{figure}
